\documentclass[10pt,a4paper]{amsart}
\usepackage[margin=19mm]{geometry}
\usepackage{amsmath,amssymb,mathtools}
\usepackage[scr=rsfs,cal=euler]{mathalfa}
\usepackage{xfrac}
\usepackage[unicode,hidelinks,bookmarks=false]{hyperref}
\newcommand{\ie}{i.e.\ }
\newcommand{\cf}{cf.\ }
\newcommand{\resp}{respectively\ }
\newcommand{\Subsection}{\subsection}
\providecommand{\nobreakdash}{\nobreak\hbox{-}}
\setlength{\emergencystretch}{2em}
\multlinegap0pt
\usepackage{amssymb}
\usepackage{float}
\usepackage{tikz}
\newtheorem{theorem}{Theorem}
\newcommand{\R}{\mathbb{R}}
\newcommand{\Z}{\mathbb{Z}}
\title{On Sum Ranges for $3n$-convergence}
\author{Preston Martens}
\date{Source: arXiv:2609.01926v1 (CC BY 4.0)}
\begin{document}
\maketitle

\noindent\emph{Source and licence.} On Sum Ranges for $3n$-convergence. Version arXiv:2609.01926v1 (CC BY 4.0), distributed by arXiv under the Creative Commons Attribution 4.0 International licence, http://creativecommons.org/licenses/by/4.0/, as recorded in the arXiv licence metadata for this exact version and shown on the abstract page https://arxiv.org/abs/2609.01926v1. Full licence notice: \texttt{licenses/arxiv-2609.01926-CC-BY-4.0.txt}.\par\smallskip

\noindent\textit{Editorial scope.} Counted unit: the article's theorem theorem (source0.tex lines 290--292). The article's complete original proof of it is delivered with it (source0.tex lines 293--299, one \texttt{proof} environment of 7 lines). No mathematics is written, repaired or re-derived: the statement and the proof are the article's own lines, in the article's own notation, and no retained line is substituted or re-flowed.  No further source-local context is needed: every cross-reference of every retained line resolves inside the two delivered spans.  Nothing was omitted from any retained span, so the package carries no omission record.  No retained line contains a citation, so the wrapper carries no bibliography and no bibliography entry.  No retained line contains a figure environment, an image-inclusion command or drawing code, so no image is delivered and the package declares no asset dependency.  Editorial formatting only, in addition: an AMS \texttt{amsart} preamble that restates the article's own declarations and macros used by the retained lines, namely the environments \texttt{theorem} on separate counters, together with the standard packages those lines need.  The retained environments print in this document as Theorem 1 in order of appearance, whereas the article prints the same items with the numbering of its own full document.  The complete native source of the article as distributed in the arXiv e-print for this exact version is bundled unchanged as \texttt{sources/209153/source0.tex}.
\begin{theorem}
    We can create infinitely many $m$-complementary sets of integers. Additionally, no two sets of zero triples share more than one element.
\end{theorem}

\begin{proof}
    We show that we can make infinitely many $2$-complement sets of integers, as higher $m$ will follow a similar process. Take some arbitrary zero triple of distinct integers $a_0 + a_1 + a_2 = 0$. To give $a_0$ a second complementary pair, we can choose $a_3 \neq a_4 \in \Z \setminus \{a_0,a_1,a_2\}$ such that $a_0 + a_3 + a_4 = 0$. Let $A_0 = \{a_0,\dots,a_4\}$. By the Pigeonhole Principle, no other triple in $A_0$ can sum to zero. Since $A_0$ is finite, the set $\{|x| : x \in A_0\}$ is bounded above by some $M_0 \in \R$. To give $a_1$ a second complement, we can choose $a_5 \neq a_6 \in \Z \setminus A_0$ such that $a_1 + a_5 + a_6 = 0$ and $|a_5|,|a_6| > 2M_0$. Since $a_5$ and $a_6$ are large, they cannot be repeats. Additionally, neither of them can complement any other existing element in $A_1 = A_0 \cup \{a_5,a_6\}$, thus ensuring we have exactly two complements for $a_0$ and $a_1$. We can continue on in this way, finding for each $a_k$ a complementary pair:
    
    \[a_{2k+3},a_{2k+4} \in \Z \setminus A_{k-1} \textnormal{ such that } |a_{2k+3}|,|a_{2k+4}| > 2M_{k-1},\]
    
    where $A_{k-1} = A_{k-2} \cup \{a_{2k+1},a_{2k+2}\}$ and $M_{k-1} \geq \max\{|x| : x \in A_{k-1}\}$. Since $a_0,a_1,a_2$ were arbitrary, we can construct a branching 2-complement set for any given zero triple.
\end{proof}

\end{document}
